% =========================================================
% capitulos/02_marco_teorico.tex
% Placeholder con las herramientas que se van a usar en la
% tesis: tablas con el estilo azul, pseudocódigo y ecuaciones
% numeradas / sin numerar de distintos tipos.
% =========================================================

\chapter{Marco Teórico}
\label{chap:marco-teorico}

Este capítulo se usará más adelante para el desarrollo teórico completo.
Por ahora contiene ejemplos de referencia de las tres herramientas que se
usarán a lo largo de la tesis: tablas, pseudocódigo de algoritmos y
ecuaciones.

\section{Ejemplo de tabla}
\label{sec:ejemplo-tabla}

La Tabla~\ref{tab:escenarios} resume los cinco escenarios considerados:

\begin{table}[htbp]
	\centering
	\label{tab:escenarios}
	
	\begin{tabular}{|c|c|c|c|c|c|}
		\hline
		\encabezado{Caso} &
		\encabezado{Muestras} &
		\encabezado{Atributos} &
		\encabezado{\begin{tabular}{c}
				Atributos\\
				expandido
		\end{tabular}} &
		\encabezado{Clases} &
		\encabezado{\begin{tabular}{c}
				Muestras\\
				por clase
		\end{tabular}} \\
		\hline
		
		1 & 150 & 3  & 9   & 3 & 50 \\
		\hline
		2 & 150 & 10 & 65  & 3 & 50 \\
		\hline
		3 & 15  & 3  & 9   & 3 & 5  \\
		\hline
		4 & 15  & 10 & 65  & 3 & 5  \\
		\hline
		5 & 15  & 15 & 135 & 3 & 5  \\
		\hline
	\end{tabular}
	\caption{Resumen de los escenarios de prueba sintéticos.}	
\end{table}

Este es el estilo estándar para \textbf{todas} las tablas de la tesis: basta
con escribir \verb|\encabezadotabla| una sola vez justo antes de la primera
celda de la fila de encabezado (después de un \verb|\hline|), y el resto de
la tabla (cuerpo, líneas, columnas) se escribe exactamente igual que una
tabla normal de LaTeX.

\section{Ejemplo de pseudocódigo}
\label{sec:ejemplo-pseudocodigo}

El Algoritmo~\ref{alg:genetico} muestra la estructura general de un
algoritmo genético, como referencia para los algoritmos que se describirán
en el desarrollo de la metodología.

\begin{algorithm}[htbp]
    \caption{Estructura general de un algoritmo genético}
    \label{alg:genetico}
    \KwEntrada{Población inicial $P_0$ de tamaño $n$, número máximo de generaciones $g_{max}$}
    \KwSalida{Mejor individuo encontrado $x^*$}

    Inicializar la población $P_0$ de forma aleatoria\;
    Evaluar la aptitud (\textit{fitness}) de cada individuo en $P_0$\;
    $g \leftarrow 0$\;

    \While{$g < g_{max}$}{
        Seleccionar individuos padres de $P_g$ según su aptitud\;
        Aplicar cruza (\textit{crossover}) para generar nuevos individuos\;
        Aplicar mutación con probabilidad $p_m$\;
        Evaluar la aptitud de los nuevos individuos\;
        $P_{g+1} \leftarrow$ nueva generación reemplazando a $P_g$\;
        $g \leftarrow g + 1$\;
    }

    $x^* \leftarrow$ individuo con mejor aptitud encontrado en todas las generaciones\;
    \Return{$x^*$}
\end{algorithm}

\section{Ejemplos de ecuaciones}
\label{sec:ejemplo-ecuaciones}

A continuación se muestran distintos formatos de ecuaciones que se usarán
en la tesis.

\subsection{Ecuación numerada simple, con lista de variables}

La forma más común: una ecuación de una línea, numerada, seguida de la
explicación de cada término.

\begin{equation}
    y = ax + b
    \label{eq:lineal-simple}
\end{equation}

\noindent donde:
\begin{itemize}
    \item $y$ es la variable dependiente (la salida del modelo),
    \item $x$ es la variable independiente (la entrada),
    \item $a$ es la pendiente de la recta,
    \item $b$ es la ordenada al origen.
\end{itemize}

\subsection{Ecuación sin numerar}

Cuando la ecuación es un paso intermedio que no necesita citarse después,
se usa el entorno sin número:

\begin{equation*}
    d(\mathbf{p}, \mathbf{q}) = \sqrt{\sum_{i=1}^{n} (p_i - q_i)^2}
\end{equation*}

que corresponde a la distancia euclidiana entre dos puntos $\mathbf{p}$ y
$\mathbf{q}$ en $\mathbb{R}^n$.

\subsection{Ecuación numerada con fracción (caso LDA/Fisher)}

Este es el tipo de ecuación que probablemente más se repita en la tesis,
al tratarse de un tema de discriminante lineal:

\begin{equation}
    J(\mathbf{w}) = \frac{\mathbf{w}^T S_B \, \mathbf{w}}{\mathbf{w}^T S_W \, \mathbf{w}}
    \label{eq:fisher-ratio}
\end{equation}

\noindent donde $S_B$ es la matriz de dispersión entre clases, $S_W$ es la
matriz de dispersión intra-clase, y $\mathbf{w}$ es el vector de proyección
que se busca optimizar.

\subsection{Sistema de ecuaciones alineado}

Cuando se necesita mostrar varios pasos o ecuaciones relacionadas, alineadas
por el signo igual, se usa \texttt{align} en vez de \texttt{equation}:

\begin{align}
    \mu_k &= \frac{1}{n_k} \sum_{i \in C_k} x_i \label{eq:media-clase}\\
    S_W &= \sum_{k=1}^{K} \sum_{i \in C_k} (x_i - \mu_k)(x_i - \mu_k)^T \label{eq:dispersion-intra}
\end{align}

\noindent la Ecuación~\eqref{eq:media-clase} calcula la media de cada clase
$C_k$, y la Ecuación~\eqref{eq:dispersion-intra} usa esas medias para
construir la matriz de dispersión intra-clase.

\subsection{Ecuación en línea (\textit{inline})}

Por último, cuando una expresión matemática es corta y forma parte de una
oración, no necesita entorno de ecuación ni número: simplemente se escribe
entre signos de peso, por ejemplo al decir que el vector de proyección
cumple $\|\mathbf{w}\| = 1$, o que el problema se resuelve para
$\mathbf{w} \in \mathbb{R}^d$.

\subsection{Formas de insertar una imagen en el documento.}

Puede insertarse una imagen sola como se muestra a continuación

\begin{figure}[htbp]
	\centering
	\includegraphics[width=0.5\textwidth]{PC1.png}
	\caption{Dimensionality Reduction techniques taxonomy}
	\label{fig:PCA BASICO}
\end{figure}


Tambien pueden insertarse varios, solo es necesario tener cuidado en como se manejan las formas, los subgráficos, los tamaños y el pie de imagen compuesto.

\begin{figure}[htbp]
	\centering
	
	% Primera fila
	\begin{subfigure}[b]{0.48\textwidth}
		\centering
		\includegraphics[width=\textwidth]{PC1.png}
		\caption{R1}
		\label{fig:R1}
	\end{subfigure}
	\hfill
	\begin{subfigure}[b]{0.48\textwidth}
		\centering
		\includegraphics[width=\textwidth]{PC2.png}
		\caption{R2}
		\label{fig:R2}
	\end{subfigure}
	
	\vspace{0.5cm}
	
	% Segunda fila
	\begin{subfigure}[b]{0.46\textwidth}
		\centering
		\includegraphics[width=\textwidth]{PC3.jpg}
		\caption{R3}
		\label{fig:R3}
	\end{subfigure}
	\hfill
	\begin{subfigure}[b]{0.45\textwidth}
		\centering
		\includegraphics[width=\textwidth]{PC4.jpg}
		\caption{R4}
		\label{fig:R4}
	\end{subfigure}
	
	\caption{Resultados obtenidos para los cuatro escenarios.}
	\label{fig:resultados_R}
\end{figure}




